\documentclass{amsart}

% 数学相关宏包
\usepackage{amsmath,mathtools,amsthm,amssymb}
\usepackage{bbm}
\usepackage{mathrsfs}
\usepackage{centernot}

% 图形和绘图
\usepackage{graphicx,float}
\usepackage{tikz,tikz-cd}
\usepackage{pgfplots}
\usepackage[framemethod=TikZ]{mdframed}
\usetikzlibrary{decorations.pathmorphing,positioning,arrows,automata,calc,shapes,patterns}
\pgfplotsset{compat=1.18}
\usepgfplotslibrary{statistics}

% 表格和列表
\usepackage{booktabs,multirow,colortbl}
\usepackage{enumitem}
\setlist[itemize,enumerate,description]{noitemsep}
\setlistdepth{9}

% 算法
\usepackage{algorithm,algorithmic}

% 字体和颜色
\usepackage{xcolor,listings}
\usepackage{xeCJK}
\setCJKmainfont{Noto Serif CJK SC}

% 页面和链接
\usepackage{fancyhdr,caption}
\usepackage{hyperref}
\usepackage{cleveref}

% 工具宏包
\usepackage{etoolbox,letltxmacro,docmute}
\usepackage{gensymb}

% 自定义设置
\renewcommand{\arraystretch}{1.5}
\let\originalmathbb\mathbb
\renewcommand{\mathbb}[1]{%
    \ifstrequal{#1}{1}{\mathbbm{#1}}{\originalmathbb{#1}}%
}

% 颜色定义
\definecolor{lightblue}{RGB}{173,216,230}
\definecolor{lightgreen}{RGB}{144,238,144}
\definecolor{lightyellow}{RGB}{255,255,224}
\definecolor{lightgray}{RGB}{211,211,211}

		% 超链接设置
		\hypersetup{
		    colorlinks=true,
		    linkcolor=red,
		    filecolor=magenta,
		    urlcolor=cyan,
		    pdftitle={\@title},
		    pdfauthor={\@author},
		    pdfsubject={数学笔记},
		    pdfkeywords={数学},
		    bookmarksnumbered=true,
		    bookmarksopen=true,
		    bookmarksopenlevel=6,
		    pdfstartview=FitH,
		    pdfpagemode=UseOutlines,
		    pdfborder={0 0 0}
		}

% 定理环境
\theoremstyle{plain}
\newtheorem{theorem}{Theorem}[section]
\newtheorem{lemma}[theorem]{Lemma}
\newtheorem{corollary}[theorem]{Corollary}
\newtheorem{proposition}[theorem]{Proposition}

\theoremstyle{definition}
\newtheorem{definition}[theorem]{Definition}
\newtheorem{example}[theorem]{Example}
\newtheorem{exercise}[theorem]{Exercise}
\newtheorem{problem}[theorem]{Problem}

\theoremstyle{remark}
\newtheorem{remark}[theorem]{Remark}
\newtheorem{note}[theorem]{Note}
\newtheorem{observation}[theorem]{Observation}

% 解答环境
\newtheorem*{solution}{Solution}
\newtheorem*{proof*}{Proof}

% 图片路径
\graphicspath{{F:/资料/2025-summer/figures/}}

\author{乐绎华, Yue Yihua}
\address{中山大学, Sun Yat-sen University}
\email{yueyh@mail2.sysu.edu.cn}
\date{\today}
\title{Mathematical Notes}

\begin{document}

\maketitle
\tableofcontents

\section{Exponential / Mixture Duality in Information Geometry}

\subsection{1. Introduction to Exponential and Mixture Families}

\subsubsection{1.1 Historical Context and Motivation}

Information geometry emerged in the 1980s through the work of Shun-ichi Amari, combining differential geometry with information theory and statistics. The field studies the \textbf{geometric structure of probability distributions}, treating families of distributions as \textbf{statistical manifolds}.

\textbf{Why is this important?}

\begin{itemize}
	\item Many fundamental problems in statistics and machine learning have elegant geometric interpretations
	\item The duality between exponential and mixture families provides a unifying framework for understanding statistical inference
	\item This geometric viewpoint leads to new algorithms and deeper theoretical insights
\end{itemize}

\subsubsection{1.2 The Fundamental Duality}

In information geometry, we study the geometric structure of probability distributions. Two fundamental types of families are \textbf{exponential families} and \textbf{mixture families}, which are dual to each other in a precise mathematical sense.

\textbf{Intuition}:

\begin{itemize}
	\item \textbf{Exponential families} are parameterized by "natural parameters" that appear linearly in the log-probability
	\item \textbf{Mixture families} are parameterized by "expectation parameters" that represent weighted averages
	\item These two parameterizations are related by the Legendre transform, creating a beautiful duality
\end{itemize}

\textbf{Key Insight}: This duality is not just mathematical convenience—it reflects a fundamental geometric structure in the space of probability distributions, analogous to how position and momentum are dual in classical mechanics through the Legendre transform.

\subsubsection{1.3 Geometric Interpretation}

Think of the space of probability distributions as a \textbf{statistical manifold}. On this manifold:

\begin{itemize}
	\item \textbf{Exponential families} form \textbf{e-flat submanifolds} (exponentially flat)
	\item \textbf{Mixture families} form \textbf{m-flat submanifolds} (mixture flat)
	\item These two types of submanifolds are \textbf{orthogonal} to each other
	\item The \textbf{Fisher information metric} provides the Riemannian structure
\end{itemize}

\textbf{Analogy}: Just as in Euclidean geometry we have horizontal and vertical lines that are orthogonal, in information geometry we have e-flat and m-flat families that are orthogonal.

\subsection{2. Exponential Families (e-families)}

\subsubsection{2.1 Definition and Basic Structure}

Given a base measure $\mu$ on a measurable space $(\Omega, \mathcal{F})$, an \textbf{exponential family} is a family of probability distributions of the form:
\[
p^{(e)}(\zeta \mid x) = \exp \left\{  \sum_{i=1}^{n} x^{i}F_i(\zeta) - \Phi(x) \right\} = \frac{1}{Z(x)}\exp \left\{  \sum_{i=1}^{n} x^{i}F_i(\zeta)  \right\}
\]

\textbf{Components Explained in Detail}:

\begin{enumerate}
	\item \textbf{Base Measure} $\mu$: The "reference" measure that doesn't depend on parameters
	\begin{itemize}
		\item For discrete distributions: counting measure
		\item For continuous distributions: Lebesgue measure
		\item Example: For Gaussian, $d\mu = d\zeta$ (Lebesgue measure)
	\end{itemize}
	\item \textbf{Natural Parameters} $x = (x^1, \ldots, x^n)$: The parameters that appear linearly in the exponent
	\begin{itemize}
		\item Also called \textbf{canonical parameters}
		\item Live in the \textbf{natural parameter space} $\mathcal{X} = \{x : \int e^{\langle x, F(\zeta) \rangle} d\mu(\zeta) < \infty\}$
		\item Form the coordinates of the \textbf{e-flat} (exponentially flat) coordinate system
	\end{itemize}
	\item \textbf{Sufficient Statistics} $F(\zeta) = [F_1(\zeta), \ldots, F_n(\zeta)]$: Fixed functions of the data
	\begin{itemize}
		\item By the \textbf{Neyman-Fisher factorization theorem}, these capture all information about the parameter
		\item The dimension $n$ is the \textbf{order} of the exponential family
		\item Example: For Gaussian, $F_1(\zeta) = \zeta$ and $F_2(\zeta) = \zeta^2$
	\end{itemize}
	\item \textbf{Log-Partition Function} $\Phi(x) = \log Z(x)$: The normalization constant
	\begin{itemize}
		\item Ensures the distribution integrates to 1
		\item Is the \textbf{cumulant generating function} of $F(\zeta)$
		\item Contains all information about the moments of the distribution
	\end{itemize}
	\item \textbf{Partition Function} $Z(x) = \int_{\Omega} \exp\left\{\sum_{i=1}^n x^i F_i(\zeta)\right\} d\mu(\zeta)$
	\begin{itemize}
		\item The normalizing constant to make it a probability distribution
		\item Must be finite for $x$ to be in the natural parameter space
	\end{itemize}
\end{enumerate}

\subsubsection{2.2 Extended Examples with Detailed Calculations}

\textbf{Example 1: Gaussian Distribution}
For $\zeta \sim \mathcal{N}(\mu, \sigma^2)$:

\textbf{Step 1: Standard form}
\[
p(\zeta) = \frac{1}{\sqrt{2\pi\sigma^2}} \exp\left(-\frac{(\zeta - \mu)^2}{2\sigma^2}\right)
\]

\textbf{Step 2: Rearrange into exponential family form}
\[
p(\zeta) = \exp\left(\frac{\mu}{\sigma^2}\zeta - \frac{1}{2\sigma^2}\zeta^2 - \frac{\mu^2}{2\sigma^2} - \frac{1}{2}\log(2\pi\sigma^2)\right)
\]

\textbf{Step 3: Identify components}

\begin{itemize}
	\item \textbf{Sufficient statistics}: $F_1(\zeta) = \zeta$, $F_2(\zeta) = \zeta^2$
	\item \textbf{Natural parameters}: $x^1 = \mu/\sigma^2$, $x^2 = -1/(2\sigma^2)$
	\item \textbf{Log-partition function}: $\Phi(x) = -\frac{(x^1)^2}{4x^2} - \frac{1}{2}\log(-2x^2) - \frac{1}{2}\log(2\pi)$
\end{itemize}

\textbf{Step 4: Verify}

\begin{itemize}
	\item $\mu = -\frac{x^1}{2x^2}$
	\item $\sigma^2 = -\frac{1}{2x^2}$
	\item Natural parameter space: $\mathcal{X} = \{(x^1, x^2) : x^2 < 0\}$
\end{itemize}

\textbf{Example 2: Categorical Distribution}
For $\zeta \in \{1, 2, \ldots, k\}$ with probabilities $p_1, \ldots, p_k$:

\textbf{Step 1: Standard form}
\[
p(\zeta = j) = p_j, \quad \sum_{j=1}^k p_j = 1
\]

\textbf{Step 2: Exponential family form}
\[
p(\zeta) = \exp\left(\sum_{j=1}^{k-1} \log\left(\frac{p_j}{p_k}\right) \mathbf{1}[\zeta = j]\right) \cdot p_k
\]

\textbf{Step 3: Identify components}

\begin{itemize}
	\item \textbf{Sufficient statistics}: $F_j(\zeta) = \mathbf{1}[\zeta = j]$ for $j = 1, \ldots, k-1$
	\item \textbf{Natural parameters}: $x^j = \log(p_j/p_k)$
	\item \textbf{Log-partition function}: $\Phi(x) = \log\left(1 + \sum_{j=1}^{k-1} e^{x^j}\right)$
\end{itemize}

\textbf{Step 4: Relationships}

\begin{itemize}
	\item $p_j = \frac{e^{x^j}}{1 + \sum_{l=1}^{k-1} e^{x^l}}$ for $j = 1, \ldots, k-1$
	\item $p_k = \frac{1}{1 + \sum_{l=1}^{k-1} e^{x^l}}$
\end{itemize}

\textbf{Example 3: Poisson Distribution}
For $\zeta \sim \text{Poisson}(\lambda)$:

\textbf{Step 1: Standard form}
\[
p(\zeta) = \frac{\lambda^\zeta e^{-\lambda}}{\zeta!}
\]

\textbf{Step 2: Exponential family form}
\[
p(\zeta) = \exp(\zeta \log \lambda - \lambda - \log \zeta!)
\]

\textbf{Step 3: Identify components}

\begin{itemize}
	\item \textbf{Sufficient statistic}: $F_1(\zeta) = \zeta$
	\item \textbf{Natural parameter}: $x^1 = \log \lambda$
	\item \textbf{Log-partition function}: $\Phi(x) = e^{x^1} = \lambda$
	\item \textbf{Base measure}: $d\mu(\zeta) = \frac{1}{\zeta!} d\zeta$ (discrete)
\end{itemize}

\subsubsection{2.3 Properties of the Log-Partition Function}

\textbf{Theorem 1}: The log-partition function $\Phi(x)$ is \textbf{strictly convex} in $x$ on the natural parameter space.

\textbf{Detailed Proof}:

\textbf{Step 1: Compute the first derivative}
\[
\frac{\partial \Phi(x)}{\partial x^i} = \frac{1}{Z(x)} \frac{\partial Z(x)}{\partial x^i} = \frac{1}{Z(x)} \int_{\Omega} F_i(\zeta) \exp\left\{\sum_{j=1}^n x^j F_j(\zeta)\right\} d\mu(\zeta)
\]

This gives us:
\[
\frac{\partial \Phi(x)}{\partial x^i} = \int_{\Omega} F_i(\zeta) p^{(e)}(\zeta \mid x) d\mu(\zeta) = \mathbb{E}_{p^{(e)}(\cdot \mid x)}[F_i(\zeta)]
\]

\textbf{Step 2: Compute the second derivative}
\[
\frac{\partial^2 \Phi(x)}{\partial x^i \partial x^j} = \frac{\partial}{\partial x^j} \mathbb{E}_{p^{(e)}(\cdot \mid x)}[F_i(\zeta)]
\]

Using the product rule and chain rule:
\[
\frac{\partial^2 \Phi(x)}{\partial x^i \partial x^j} = \mathbb{E}_{p^{(e)}(\cdot \mid x)}[F_i(\zeta) F_j(\zeta)] - \mathbb{E}_{p^{(e)}(\cdot \mid x)}[F_i(\zeta)] \mathbb{E}_{p^{(e)}(\cdot \mid x)}[F_j(\zeta)]
\]

\textbf{Step 3: Recognize as covariance}
\[
\frac{\partial^2 \Phi(x)}{\partial x^i \partial x^j} = \text{Cov}_{p^{(e)}(\cdot \mid x)}[F_i(\zeta), F_j(\zeta)]
\]

\textbf{Step 4: Conclude convexity}
Since covariance matrices are positive semidefinite, the Hessian matrix $\nabla^2 \Phi(x)$ is positive semidefinite, making $\Phi(x)$ convex. It's strictly convex if the covariance matrix is positive definite, which occurs when the sufficient statistics are not linearly dependent.

\textbf{Theorem 2}: The Hessian of $\Phi(x)$ equals the Fisher information matrix.

\textbf{Proof}: The Fisher information matrix is defined as:
\[
\mathcal{I}_{ij}(x) = \mathbb{E}_{p^{(e)}(\cdot \mid x)}\left[\frac{\partial \log p^{(e)}(\zeta \mid x)}{\partial x^i} \frac{\partial \log p^{(e)}(\zeta \mid x)}{\partial x^j}\right]
\]

Since $\frac{\partial \log p^{(e)}(\zeta \mid x)}{\partial x^i} = F_i(\zeta) - \frac{\partial \Phi(x)}{\partial x^i}$, we have:
\[
\mathcal{I}_{ij}(x) = \mathbb{E}_{p^{(e)}(\cdot \mid x)}\left[\left(F_i(\zeta) - \mathbb{E}[F_i]\right)\left(F_j(\zeta) - \mathbb{E}[F_j]\right)\right] = \frac{\partial^2 \Phi(x)}{\partial x^i \partial x^j}
\]

\textbf{Corollary}: The Fisher information metric on the exponential family manifold is given by:
\[
g_{ij}(x) = \frac{\partial^2 \Phi(x)}{\partial x^i \partial x^j}
\]

\subsubsection{2.4 Moment Generating Properties}

\textbf{Theorem 3}: The log-partition function is the \textbf{cumulant generating function} of the sufficient statistics.

\textbf{Proof}: The moment generating function of $F(\zeta)$ under $p^{(e)}(\cdot \mid x)$ is:
\[
M_F(t) = \mathbb{E}_{p^{(e)}(\cdot \mid x)}[e^{\langle t, F(\zeta) \rangle}] = \frac{Z(x + t)}{Z(x)} = e^{\Phi(x + t) - \Phi(x)}
\]

The cumulant generating function is:
\[
K_F(t) = \log M_F(t) = \Phi(x + t) - \Phi(x)
\]

\textbf{Consequences}:

\begin{enumerate}
	\item \textbf{First moment}: $\mathbb{E}[F_i] = \frac{\partial \Phi(x)}{\partial x^i}$
	\item \textbf{Second moment}: $\mathbb{E}[F_i F_j] = \frac{\partial^2 \Phi(x)}{\partial x^i \partial x^j} + \frac{\partial \Phi(x)}{\partial x^i} \frac{\partial \Phi(x)}{\partial x^j}$
	\item \textbf{Higher moments}: Can be computed by taking higher derivatives of $\Phi(x)$
\end{enumerate}

\subsubsection{2.5 Maximum Likelihood Estimation}

\textbf{Theorem 4}: For observed data $\zeta$, the maximum likelihood estimate (MLE) of $x$ is:
\[
\widehat{x}(\zeta) = \arg\max_x \log p^{(e)}(\zeta \mid x) = \arg\max_x \left( \sum_{i}x^{i}F_i(\zeta) - \Phi(x) \right)
\]

\textbf{Detailed Solution}:

\textbf{Step 1: Set up the optimization problem}
\[
\max_x \ell(x) = \max_x \left( \sum_{i=1}^n x^{i}F_i(\zeta) - \Phi(x) \right)
\]

\textbf{Step 2: Take the first-order condition}
\[
\frac{\partial \ell(x)}{\partial x^i} = F_i(\zeta) - \frac{\partial \Phi(x)}{\partial x^i} = 0
\]

\textbf{Step 3: Solve for the MLE}
\[
F_i(\zeta) = \frac{\partial \Phi(\widehat{x})}{\partial x^i} = \mathbb{E}_{p^{(e)}(\cdot \mid \widehat{x})}[F_i(\zeta)]
\]

\textbf{Interpretation}: The MLE equates the observed sufficient statistics with their expected values under the fitted distribution.

\textbf{Step 4: Express in terms of conjugate function}
The maximum value of the log-likelihood is:
\[
\max_x \ell(x) = \Phi^{*}(F(\zeta))
\]

where $\Phi^{*}$ is the \textbf{conjugate function} (Legendre transform) of $\Phi$:
\[
\Phi^{*}(u) = \sup_{x} \left\{ \langle x, u \rangle - \Phi(x) \right\}
\]

\textbf{Geometric Interpretation}: The MLE finds the point $x$ where the hyperplane $\langle x, F(\zeta) \rangle$ is tangent to the convex function $\Phi(x)$. This is the essence of the Legendre transform.

\subsubsection{2.6 Exponential Family Manifold Structure}

\textbf{Coordinate Charts}: An exponential family can be viewed as a \textbf{statistical manifold} with two natural coordinate systems:

\begin{enumerate}
	\item \textbf{Natural coordinates} $(x^1, \ldots, x^n)$: The natural parameters
	\item \textbf{Expectation coordinates} $(u_1, \ldots, u_n)$: The expectation parameters
\end{enumerate}

\textbf{Metric Structure}: The Fisher information provides a Riemannian metric:
\[
ds^2 = \sum_{i,j} g_{ij}(x) dx^i dx^j = \sum_{i,j} \frac{\partial^2 \Phi(x)}{\partial x^i \partial x^j} dx^i dx^j
\]

\textbf{Geodesics}: In natural coordinates, the geodesics are straight lines. This is why exponential families are called "exponentially flat" or "e-flat".

\textbf{Connection}: The \textbf{Levi-Civita connection} in natural coordinates has all \textbf{Christoffel symbols} equal to zero:
\[
\Gamma_{ij}^k = 0 \quad \text{for all } i, j, k
\]

This means the manifold is \textbf{flat} in the natural coordinate system.

\subsection{3. The Duality Relationship}

\subsubsection{3.1 Dual Parameterizations and the Legendre Transform}

The heart of exponential/mixture duality lies in the \textbf{Legendre transform}, which provides a bijection between the natural parameter space and the expectation parameter space.

\textbf{Definition}: For a convex function $\Phi(x)$, its \textbf{Legendre transform} (or \textbf{convex conjugate}) is:
\[
\Phi^{*}(u) = \sup_{x} \left\{ \langle x, u \rangle - \Phi(x) \right\}
\]

\textbf{Geometric Interpretation}:

\begin{itemize}
	\item $\Phi^{*}(u)$ is the maximum "height" of the hyperplane $\langle x, u \rangle$ above the graph of $\Phi(x)$
	\item The supremum is achieved when the hyperplane is tangent to the graph of $\Phi(x)$
	\item This creates a duality between points $x$ and hyperplanes with slope $u$
\end{itemize}

\subsubsection{3.2 Detailed Derivation of Duality Relations}

\textbf{Step 1: Start with the MLE condition}
At the MLE $\widehat{x}(\zeta)$, we have:
\[
\log p^{(e)}(\zeta \mid \widehat{x}(\zeta)) = \langle \widehat{x}(\zeta), F(\zeta) \rangle - \Phi(\widehat{x}(\zeta)) = \Phi^{*}(F(\zeta))
\]

\textbf{Step 2: Define expectation parameters}
Let $u = F(\zeta)$ be the \textbf{expectation parameters}. The key insight is that these represent the expected values of the sufficient statistics:
\[
u_i = \mathbb{E}_{p^{(e)}(\cdot \mid x)}[F_i(\zeta)] = \frac{\partial \Phi(x)}{\partial x^i}
\]

\textbf{Step 3: Establish the duality}
From the definition of the Legendre transform and the first-order condition:
\[
\Phi(\widehat{x}) - \langle \widehat{x}, u \rangle + \Phi^{*}(u) = 0
\]

This gives us the \textbf{fundamental duality relations}:
\[
u_i = \frac{\partial \Phi(x)}{\partial x^{i}} \quad \text{and} \quad x^{i} = \frac{\partial \Phi^{*}(u)}{\partial u_i}
\]

\textbf{Step 4: Prove the inverse relationship}
Since $\Phi$ is strictly convex, the mapping $x \mapsto \nabla \Phi(x)$ is bijective. Moreover:
\[
\Phi^{**}(x) = \Phi(x) \quad \text{(involution property)}
\]

This means the Legendre transform is its own inverse when applied twice.

\subsubsection{3.3 Bregman Divergences and their Geometric Meaning}

\textbf{Definition}: The \textbf{Bregman divergence} associated with a convex function $\Phi$ is:
\[
B_{\Phi}(x, y) = \Phi(x) - \Phi(y) - \langle \nabla \Phi(y), x - y \rangle
\]

\textbf{Geometric Interpretation}:

\begin{itemize}
	\item $B_{\Phi}(x, y)$ measures the "distance" from the point $x$ to the tangent hyperplane of $\Phi$ at point $y$
	\item It's the difference between the actual function value $\Phi(x)$ and its first-order Taylor approximation around $y$
	\item Always non-negative due to convexity of $\Phi$
\end{itemize}

\textbf{Properties}:

\begin{enumerate}
	\item \textbf{Non-negativity}: $B_{\Phi}(x, y) \geq 0$ with equality iff $x = y$
	\item \textbf{Asymmetry}: $B_{\Phi}(x, y) \neq B_{\Phi}(y, x)$ in general
	\item \textbf{Convexity}: $B_{\Phi}(x, y)$ is convex in $x$ for fixed $y$
\end{enumerate}

\textbf{Connection to Log-Likelihood}:
The log-likelihood can be expressed as:
\[
\log p^{(e)}(\zeta \mid x) = -B_{\Phi}(x, \widehat{x}(\zeta)) + \log p^{(e)}(\zeta \mid \widehat{x}(\zeta))
\]

where $\widehat{x}(\zeta)$ is the MLE for data $\zeta$.

\subsubsection{3.4 KL Divergence as Bregman Divergence}

\textbf{Theorem}: The KL divergence between two exponential family distributions is a Bregman divergence:
\[
\mathbb{K}[p^{(e)}(\cdot \mid x') \parallel p^{(e)}(\cdot \mid x)] = B_{\Phi}(x', x)
\]

\textbf{Detailed Proof}:

\textbf{Step 1: Define the KL divergence}
\[
\mathbb{K}[p^{(e)}(\cdot \mid x') \parallel p^{(e)}(\cdot \mid x)] = \int_{\Omega} p^{(e)}(\zeta \mid x') \log \frac{p^{(e)}(\zeta \mid x')}{p^{(e)}(\zeta \mid x)} d\mu(\zeta)
\]

\textbf{Step 2: Substitute the exponential family form}
\[
= \int_{\Omega} p^{(e)}(\zeta \mid x') \left[ \sum_{i=1}^n (x'^i - x^i) F_i(\zeta) - \Phi(x') + \Phi(x) \right] d\mu(\zeta)
\]

\textbf{Step 3: Use linearity of expectation}
\[
= \sum_{i=1}^n (x'^i - x^i) \mathbb{E}_{p^{(e)}(\cdot \mid x')}[F_i(\zeta)] - \Phi(x') + \Phi(x)
\]

\textbf{Step 4: Apply the moment condition}
Since $\mathbb{E}_{p^{(e)}(\cdot \mid x')}[F_i(\zeta)] = \frac{\partial \Phi(x')}{\partial x'^i}$:
\[
= \sum_{i=1}^n (x'^i - x^i) \frac{\partial \Phi(x')}{\partial x'^i} - \Phi(x') + \Phi(x)
\]

\textbf{Step 5: Recognize as Bregman divergence}
\[
= \langle x' - x, \nabla \Phi(x') \rangle - \Phi(x') + \Phi(x) = B_{\Phi}(x', x)
\]

\textbf{Corollary}: This establishes that the KL divergence is the natural "distance" measure for exponential families, and it has all the properties of Bregman divergences.

\subsubsection{3.5 Dual Coordinate Systems}

\textbf{Natural Coordinates} $(x^1, \ldots, x^n)$:

\begin{itemize}
	\item Parameters that appear linearly in the log-density
	\item Form the \textbf{e-flat} (exponentially flat) coordinate system
	\item Geodesics are straight lines in this coordinate system
	\item The Fisher metric in these coordinates is $g_{ij} = \frac{\partial^2 \Phi}{\partial x^i \partial x^j}$
\end{itemize}

\textbf{Expectation Coordinates} $(u_1, \ldots, u_n)$:

\begin{itemize}
	\item Expected values of the sufficient statistics
	\item Form the \textbf{m-flat} (mixture flat) coordinate system
	\item Mixture curves are straight lines in this coordinate system
	\item The Fisher metric in these coordinates is $g^{ij} = \frac{\partial^2 \Phi^*}{\partial u_i \partial u_j}$
\end{itemize}

\textbf{Transformation Between Coordinates}:

\begin{itemize}
	\item \textbf{Forward}: $u_i = \frac{\partial \Phi(x)}{\partial x^i}$
	\item \textbf{Inverse}: $x^i = \frac{\partial \Phi^*(u)}{\partial u_i}$
\end{itemize}

\textbf{Metric Relationship}:
The metrics in the two coordinate systems are inverse to each other:
\[
g^{ij}(u) = \left[g_{kl}(x)\right]^{-1} \frac{\partial x^k}{\partial u_i} \frac{\partial x^l}{\partial u_j}
\]

\subsection{4. Mixture Families (m-families)}

\subsubsection{4.1 Definition and Structure}

A \textbf{mixture family} consists of distributions formed by taking convex combinations of base distributions:
\[
p^{(m)}(\zeta \mid \eta) = \sum_{j=1}^k \eta_j q_j(\zeta)
\]

where:

\begin{itemize}
	\item $\eta = (\eta_1, \ldots, \eta_k)$ are the \textbf{mixing weights}
	\item $\eta_j \geq 0$ and $\sum_{j=1}^k \eta_j = 1$ (simplex constraint)
	\item $q_j(\zeta)$ are the \textbf{base distributions} or \textbf{component distributions}
\end{itemize}

\textbf{Geometric Structure}:

\begin{itemize}
	\item The mixture family forms a \textbf{convex hull} of the base distributions
	\item The parameter space is the \textbf{probability simplex} $\Delta_{k-1}$
	\item The family is \textbf{m-flat} (mixture flat) - convex combinations are straight lines
\end{itemize}

\subsubsection{4.2 Exponential vs Mixture Curves}

\textbf{m-mixture} (mixture family curve):
For two distributions $p(\zeta)$ and $q(\zeta)$, the m-mixture is:
\[
r_t(\zeta) = tp(\zeta) + (1-t)q(\zeta), \quad t \in [0,1]
\]

\textbf{Properties}:

\begin{itemize}
	\item This is a \textbf{straight line} in the space of probability distributions
	\item The curve lies in the \textbf{convex hull} of $p$ and $q$
	\item The parameterization is \textbf{affine} in the mixing parameter $t$
\end{itemize}

\textbf{e-mixture} (exponential family curve):
For two distributions $p(\zeta)$ and $q(\zeta)$, the e-mixture is:
\[
r_t(\zeta) = \frac{p(\zeta)^t q(\zeta)^{1-t}}{Z_t}, \quad t \in [0,1]
\]

where $Z_t = \int_{\Omega} p(\zeta)^t q(\zeta)^{1-t} d\mu(\zeta)$ is the normalization constant.

\textbf{Properties}:

\begin{itemize}
	\item This is a \textbf{geodesic} in the exponential family manifold
	\item The curve follows the \textbf{geometric mean} structure
	\item The parameterization is \textbf{exponential} in the mixing parameter $t$
\end{itemize}

\subsubsection{4.3 Detailed Geometric Analysis}

\textbf{Orthogonality}:
The e-flat and m-flat curves are \textbf{orthogonal} in the sense of the Fisher information metric. This means:
\[
\langle \dot{\gamma}_e(t), \dot{\gamma}_m(t) \rangle_{\text{Fisher}} = 0
\]

where $\gamma_e(t)$ and $\gamma_m(t)$ are e-flat and m-flat curves intersecting at a point.

\textbf{Pythagorean Theorem}:
For three distributions $p$, $q$, $r$ where $q$ is the orthogonal projection of $p$ onto the m-flat submanifold containing $r$:
\[
\mathbb{K}[p \parallel r] = \mathbb{K}[p \parallel q] + \mathbb{K}[q \parallel r]
\]

\textbf{Geometric Interpretation}:

\begin{itemize}
	\item The "distance" (KL divergence) from $p$ to $r$ decomposes into two orthogonal components
	\item One component is the "distance" from $p$ to its projection $q$
	\item The other component is the "distance" from the projection $q$ to $r$
\end{itemize}

\subsubsection{4.4 Dual Coordinate Representation}

\textbf{In Natural Coordinates}:
If the base distributions $q_j(\zeta)$ are from the same exponential family:
\[
q_j(\zeta) = \exp\left\{\sum_{i=1}^n x_j^i F_i(\zeta) - \Phi(x_j)\right\}
\]

Then the mixture can be written as:
\[
p^{(m)}(\zeta \mid \eta) = \sum_{j=1}^k \eta_j \exp\left\{\sum_{i=1}^n x_j^i F_i(\zeta) - \Phi(x_j)\right\}
\]

\textbf{In Expectation Coordinates}:
The mixture has expectation parameters:
\[
u_i = \sum_{j=1}^k \eta_j u_{j,i}
\]

where $u_{j,i}$ are the expectation parameters of the $j$-th component.

\textbf{Key Insight}: Mixture families are \textbf{convex} in expectation coordinates but \textbf{non-convex} in natural coordinates.

\subsubsection{4.5 Applications of Mixture Families}

\textbf{Finite Mixture Models}:

\begin{itemize}
	\item \textbf{Gaussian Mixture Models}: $p(\zeta) = \sum_{j=1}^k \pi_j \mathcal{N}(\zeta; \mu_j, \Sigma_j)$
	\item \textbf{Mixture of Experts}: Each component specializes in different regions
	\item \textbf{Clustering}: Each component represents a different cluster
\end{itemize}

\textbf{Hierarchical Models}:

\begin{itemize}
	\item \textbf{Bayesian Inference}: Prior distributions as mixtures
	\item \textbf{Empirical Bayes}: Estimating the mixing distribution
	\item \textbf{Nonparametric Methods}: Infinite mixture models (Dirichlet processes)
\end{itemize}

\textbf{Approximation Theory}:

\begin{itemize}
	\item \textbf{Density Estimation}: Approximating complex distributions
	\item \textbf{Variational Inference}: Approximating posterior distributions
	\item \textbf{Neural Networks}: Mixture density networks
\end{itemize}

\subsubsection{4.6 Dual Representation Theorem}

\textbf{Theorem}: Every mixture family can be represented as an exponential family in a higher-dimensional space, and vice versa.

\textbf{Proof Sketch}:

\textbf{Step 1: Mixture to Exponential}
Consider a mixture family:
\[
p^{(m)}(\zeta \mid \eta) = \sum_{j=1}^k \eta_j q_j(\zeta)
\]

Introduce latent variables $z \in \{1, \ldots, k\}$ and define:
\[
p(\zeta, z \mid \eta) = \eta_z q_z(\zeta)
\]

This is an exponential family in the joint space $(\zeta, z)$.

\textbf{Step 2: Exponential to Mixture}
Consider an exponential family:
\[
p^{(e)}(\zeta \mid x) = \exp\left\{\sum_{i=1}^n x^i F_i(\zeta) - \Phi(x)\right\}
\]

This can be written as a mixture over the sufficient statistic space:
\[
p^{(e)}(\zeta \mid x) = \int p(\zeta \mid u) p(u \mid x) du
\]

where $p(u \mid x)$ is the distribution of the sufficient statistics.

\textbf{Implications}:

\begin{itemize}
	\item The exponential/mixture duality is not just a mathematical curiosity
	\item It provides practical algorithms for inference and learning
	\item EM algorithm exploits this duality for mixture model fitting
\end{itemize}

\subsection{5. Information-Theoretic Quantities}

\subsubsection{5.1 Shannon Entropy}

For an exponential family parameterized by expectation parameters $u$:
\[
\mathbb{S}[p(\cdot \mid x(u))] = S(u) = \Phi^{*}(u)
\]

\textbf{Detailed Derivation}:
\[
\mathbb{S}[p] = -\int_{\Omega} p(\zeta) \log p(\zeta) d\mu(\zeta)
\]

For an exponential family $p^{(e)}(\zeta \mid x)$:
\[
\mathbb{S}[p^{(e)}] = -\int_{\Omega} p^{(e)}(\zeta \mid x) \left[\sum_{i=1}^n x^i F_i(\zeta) - \Phi(x)\right] d\mu(\zeta)
\]

Using the expectation parameter relationship $u_i = \mathbb{E}[F_i(\zeta)]$:
\[
\mathbb{S}[p^{(e)}] = -\sum_{i=1}^n x^i u_i + \Phi(x)
\]

From the duality relation $x^i = \frac{\partial \Phi^*(u)}{\partial u_i}$:
\[
\mathbb{S}[p^{(e)}] = -\sum_{i=1}^n \frac{\partial \Phi^*(u)}{\partial u_i} u_i + \Phi\left(\frac{\partial \Phi^*(u)}{\partial u}\right)
\]

By the \textbf{Euler's theorem} for homogeneous functions and properties of the Legendre transform:
\[
\mathbb{S}[p^{(e)}] = \Phi^{*}(u)
\]

\textbf{Geometric Interpretation}: The entropy is the "height" of the tangent hyperplane to $\Phi(x)$ at the point corresponding to expectation parameters $u$.

\subsubsection{5.2 Cross-Entropy and Mutual Information}

\textbf{Cross-Entropy}:
The cross-entropy between a distribution $p(\zeta)$ and an exponential family $p^{(e)}(\zeta \mid x)$ is:
\[
H(p, p^{(e)}) = -\int_{\Omega} p(\zeta) \log p^{(e)}(\zeta \mid x) d\mu(\zeta)
\]

\textbf{Detailed Calculation}:
\[
H(p, p^{(e)}) = -\int_{\Omega} p(\zeta) \left[\sum_{i=1}^n x^i F_i(\zeta) - \Phi(x)\right] d\mu(\zeta)
\]
\[
= -\sum_{i=1}^n x^i \mathbb{E}_p[F_i(\zeta)] + \Phi(x)
\]
\[
= \Phi(x) - \langle x, \mathbb{E}_p[F(\zeta)] \rangle
\]

\textbf{Mutual Information}:
For a joint exponential family $p(\zeta, \xi \mid x)$, the mutual information is:
\[
I(\zeta; \xi) = \mathbb{K}[p(\zeta, \xi) \parallel p(\zeta)p(\xi)]
\]

This can be expressed as a difference of log-partition functions:
\[
I(\zeta; \xi) = \Phi_{joint}(x) - \Phi_{\zeta}(x_{\zeta}) - \Phi_{\xi}(x_{\xi})
\]

\subsubsection{5.3 Relative Entropy and Information Divergence}

\textbf{Relative Entropy} (KL divergence):
\[
\mathbb{K}[p \parallel q] = \int_{\Omega} p(\zeta) \log \frac{p(\zeta)}{q(\zeta)} d\mu(\zeta)
\]

\textbf{Information Divergence Properties}:

\begin{enumerate}
	\item \textbf{Data Processing Inequality}: For any function $T(\zeta)$:
\[
\mathbb{K}[p(T(\zeta)) \parallel q(T(\zeta))] \leq \mathbb{K}[p(\zeta) \parallel q(\zeta)]
\]
	\item \textbf{Chain Rule}: For joint distributions:
\[
\mathbb{K}[p(\zeta, \xi) \parallel q(\zeta, \xi)] = \mathbb{K}[p(\zeta) \parallel q(\zeta)] + \mathbb{K}[p(\xi \mid \zeta) \parallel q(\xi \mid \zeta)]
\]
	\item \textbf{Convexity}: $\mathbb{K}[p \parallel q]$ is convex in both arguments
\end{enumerate}

\textbf{Applications}:

\begin{itemize}
	\item \textbf{Hypothesis Testing}: KL divergence measures discriminability between hypotheses
	\item \textbf{Compression}: Relates to optimal coding lengths
	\item \textbf{Statistical Inference}: Measures efficiency of estimators
\end{itemize}

\subsection{6. Maximum Entropy and Minimum Divergence Inference}

\subsubsection{6.1 Maximum Entropy Principle}

\textbf{Problem Statement}: Find the distribution that maximizes entropy subject to moment constraints:
\[
\max_p \mathbb{S}[p] \quad \text{subject to} \quad \mathbb{E}_p[F_i(\zeta)] = \mu_i, \quad i = 1, \ldots, n
\]

\textbf{Lagrangian Formulation}:
\[
\mathcal{L}[p, \lambda] = \mathbb{S}[p] - \sum_{i=1}^n \lambda_i \left(\mathbb{E}_p[F_i(\zeta)] - \mu_i\right)
\]

\textbf{Variational Solution}:
Taking the functional derivative with respect to $p(\zeta)$:
\[
\frac{\delta \mathcal{L}}{\delta p} = -\log p(\zeta) - 1 - \sum_{i=1}^n \lambda_i F_i(\zeta) = 0
\]

This gives:
\[
p^{*}(\zeta) = \exp\left\{-1 - \sum_{i=1}^n \lambda_i F_i(\zeta)\right\}
\]

\textbf{Normalization}:
\[
p^{*}(\zeta) = \exp\left\{\sum_{i=1}^n \lambda_i F_i(\zeta) - \Phi(\lambda)\right\}
\]

where $\Phi(\lambda) = \log \int_{\Omega} \exp\left\{\sum_{i=1}^n \lambda_i F_i(\zeta)\right\} d\mu(\zeta)$.

\textbf{Key Result}: The maximum entropy distribution is an exponential family with natural parameters $\lambda$.

\subsubsection{6.2 Detailed Examples of Maximum Entropy}

\textbf{Example 1: Uniform Distribution}

\begin{itemize}
	\item \textbf{Constraint}: None (just normalization)
	\item \textbf{Result}: $p(\zeta) = \frac{1}{|\Omega|}$ (uniform distribution)
	\item \textbf{Interpretation}: Maximum entropy with no constraints
\end{itemize}

\textbf{Example 2: Exponential Distribution}

\begin{itemize}
	\item \textbf{Constraint}: $\mathbb{E}[\zeta] = \mu$ for $\zeta \geq 0$
	\item \textbf{Sufficient statistic}: $F_1(\zeta) = \zeta$
	\item \textbf{Result}: $p(\zeta) = \frac{1}{\mu} e^{-\zeta/\mu}$
	\item \textbf{Natural parameter}: $\lambda = -1/\mu$
\end{itemize}

\textbf{Example 3: Gaussian Distribution}

\begin{itemize}
	\item \textbf{Constraints}: $\mathbb{E}[\zeta] = \mu$, $\mathbb{E}[\zeta^2] = \sigma^2 + \mu^2$
	\item \textbf{Sufficient statistics}: $F_1(\zeta) = \zeta$, $F_2(\zeta) = \zeta^2$
	\item \textbf{Result}: $p(\zeta) = \frac{1}{\sqrt{2\pi\sigma^2}} \exp\left(-\frac{(\zeta-\mu)^2}{2\sigma^2}\right)$
\end{itemize}

\subsubsection{6.3 Minimum Divergence Principle}

\textbf{Problem Statement}: Find the distribution in a given family that minimizes KL divergence to a target distribution:
\[
\min_{p \in \mathcal{P}} \mathbb{K}[p \parallel p_{\text{target}}]
\]

\textbf{For Exponential Families}:
\[
\min_x \mathbb{K}[p^{(e)}(\cdot \mid x) \parallel p_{\text{target}}]
\]

\textbf{Solution}:
Taking the derivative with respect to $x^i$:
\[
\frac{\partial}{\partial x^i} \mathbb{K}[p^{(e)}(\cdot \mid x) \parallel p_{\text{target}}] = \mathbb{E}_{p^{(e)}(\cdot \mid x)}[F_i(\zeta)] - \mathbb{E}_{p_{\text{target}}}[F_i(\zeta)] = 0
\]

This gives the \textbf{moment matching condition}:
\[
\mathbb{E}_{p^{(e)}(\cdot \mid x^*)}[F_i(\zeta)] = \mathbb{E}_{p_{\text{target}}}[F_i(\zeta)]
\]

\textbf{Geometric Interpretation}: The optimal distribution is the one whose sufficient statistics have the same expected values as the target distribution.

\subsubsection{6.4 Algorithms for Maximum Entropy}

\textbf{Iterative Scaling Algorithm}:

\begin{enumerate}
	\item Initialize $x^{(0)}$
	\item For $t = 0, 1, 2, \ldots$:
	\begin{itemize}
		\item Compute $\mathbb{E}_{p^{(e)}(\cdot \mid x^{(t)})}[F_i(\zeta)]$
		\item Update: $x^{(t+1)}_i = x^{(t)}_i + \alpha \left(\mu_i - \mathbb{E}_{p^{(e)}(\cdot \mid x^{(t)})}[F_i(\zeta)]\right)$
	\end{itemize}
	\item Repeat until convergence
\end{enumerate}

\textbf{Newton-Raphson Method}:
Using the Fisher information matrix as the Hessian:
\[
x^{(t+1)} = x^{(t)} + \mathcal{I}^{-1}(x^{(t)}) \left(\mu - \mathbb{E}_{p^{(e)}(\cdot \mid x^{(t)})}[F(\zeta)]\right)
\]

\subsection{7. Geometric Properties of KL Divergence}

\subsubsection{7.1 KL Divergence as a "Distance"}

The KL divergence $\mathbb{K}[p \parallel q]$ is \textbf{not} a true distance metric, but it has useful properties:

\begin{enumerate}
	\item \textbf{Non-negativity}: $\mathbb{K}[p \parallel q] \geq 0$, with equality iff $p = q$
	\item \textbf{Asymmetry}: $\mathbb{K}[p \parallel q] \neq \mathbb{K}[q \parallel p]$ in general
	\item \textbf{Convexity}: $\mathbb{K}[p \parallel q]$ is convex in both arguments
\end{enumerate}

\subsubsection{7.2 Generalized Pythagorean Theorem}

For probability distributions $p(\cdot)$, $q(\cdot)$, $s(\cdot)$:
\[
\mathbb{K}[p \parallel s] = \mathbb{K}[p \parallel q] + \mathbb{K}[q \parallel s] + \int_{\Omega} (p - q)(\log s - \log q) d\mu
\]

\textbf{Special Case}: If $q$ is the orthogonal projection of $p$ onto the manifold containing $s$, then:
\[
\mathbb{K}[p \parallel s] = \mathbb{K}[p \parallel q] + \mathbb{K}[q \parallel s]
\]

This is the \textbf{Pythagorean theorem} in information geometry!

\textbf{Geometric Interpretation}: The "distance" from $p$ to $s$ can be decomposed into the "distance" from $p$ to its projection $q$ plus the "distance" from $q$ to $s$.

\subsubsection{7.3 Detailed Proof of the Pythagorean Theorem}

\textbf{Setup}: Let $\mathcal{M}$ be an m-flat submanifold and $\mathcal{E}$ be an e-flat submanifold such that they intersect orthogonally at point $q$.

\textbf{Step 1}: For any $p \in \mathcal{E}$ and $s \in \mathcal{M}$:
\[
\mathbb{K}[p \parallel s] = \int_{\Omega} p(\zeta) \log \frac{p(\zeta)}{s(\zeta)} d\mu(\zeta)
\]

\textbf{Step 2}: Decompose the logarithm:
\[
\log \frac{p(\zeta)}{s(\zeta)} = \log \frac{p(\zeta)}{q(\zeta)} + \log \frac{q(\zeta)}{s(\zeta)}
\]

\textbf{Step 3}: Apply linearity:
\[
\mathbb{K}[p \parallel s] = \int_{\Omega} p(\zeta) \log \frac{p(\zeta)}{q(\zeta)} d\mu(\zeta) + \int_{\Omega} p(\zeta) \log \frac{q(\zeta)}{s(\zeta)} d\mu(\zeta)
\]

\textbf{Step 4}: Use the orthogonality condition:
The orthogonality means that $\int_{\Omega} p(\zeta) \log \frac{q(\zeta)}{s(\zeta)} d\mu(\zeta) = \int_{\Omega} q(\zeta) \log \frac{q(\zeta)}{s(\zeta)} d\mu(\zeta)$.

\textbf{Step 5}: Conclude:
\[
\mathbb{K}[p \parallel s] = \mathbb{K}[p \parallel q] + \mathbb{K}[q \parallel s]
\]

\subsubsection{7.4 Applications of the Pythagorean Theorem}

\textbf{Variational Inference}:

\begin{itemize}
	\item Find the best approximation $q^*$ in a family $\mathcal{Q}$ to a target $p$
	\item Decompose the approximation error using orthogonal projections
	\item Guides the choice of variational families
\end{itemize}

\textbf{Model Selection}:

\begin{itemize}
	\item Compare models by their projections onto different submanifolds
	\item Balance between model complexity and goodness of fit
	\item Provides geometric understanding of bias-variance tradeoff
\end{itemize}

\textbf{Information Bottleneck}:

\begin{itemize}
	\item Compress information while preserving relevant structure
	\item Use orthogonal projections to separate relevant and irrelevant information
	\item Applications in neural networks and representation learning
\end{itemize}

\subsection{8. Applications and Examples}

\subsubsection{8.1 Gaussian Distribution}

For $\zeta \sim \mathcal{N}(\mu, \sigma^2)$:

\begin{itemize}
	\item \textbf{Sufficient statistics}: $F_1(\zeta) = \zeta$, $F_2(\zeta) = \zeta^2$
	\item \textbf{Natural parameters}: $x^1 = \mu/\sigma^2$, $x^2 = -1/(2\sigma^2)$
	\item \textbf{Log-partition function}: $\Phi(x) = -\frac{(x^1)^2}{4x^2} - \frac{1}{2}\log(-2x^2) - \frac{1}{2}\log(2\pi)$
	\item \textbf{Expectation parameters}: $u_1 = \mu$, $u_2 = \mu^2 + \sigma^2$
\end{itemize}

\textbf{Multivariate Gaussian}:
For $\zeta \sim \mathcal{N}(\mu, \Sigma)$:

\begin{itemize}
	\item \textbf{Sufficient statistics}: $F_1(\zeta) = \zeta$, $F_2(\zeta) = \zeta \zeta^T$
	\item \textbf{Natural parameters}: $x_1 = \Sigma^{-1}\mu$, $X_2 = -\frac{1}{2}\Sigma^{-1}$
	\item \textbf{Log-partition function}: $\Phi(x) = \frac{1}{2}\mu^T\Sigma^{-1}\mu + \frac{1}{2}\log|\Sigma| + \frac{d}{2}\log(2\pi)$
\end{itemize}

\subsubsection{8.2 Categorical Distribution}

For $\zeta \in \{1, 2, \ldots, k\}$ with probabilities $p_1, \ldots, p_k$:

\begin{itemize}
	\item \textbf{Sufficient statistics}: $F_i(\zeta) = \mathbf{1}[\zeta = i]$
	\item \textbf{Natural parameters}: $x^i = \log p_i$ (up to normalization)
	\item \textbf{Log-partition function}: $\Phi(x) = \log\left(\sum_{i=1}^k e^{x^i}\right)$
	\item \textbf{Expectation parameters}: $u_i = p_i$
\end{itemize}

\textbf{Multinomial Distribution}:
For $n$ trials with outcome counts $\zeta = (n_1, \ldots, n_k)$:

\begin{itemize}
	\item \textbf{Sufficient statistics}: $F_i(\zeta) = n_i$
	\item \textbf{Natural parameters}: $x^i = \log p_i$
	\item \textbf{Log-partition function}: $\Phi(x) = n \log\left(\sum_{i=1}^k e^{x^i}\right)$
\end{itemize}

\subsubsection{8.3 Exponential Family GLMs}

\textbf{Generalized Linear Models} extend linear regression to exponential families:
\[
\mathbb{E}[Y \mid X] = g^{-1}(X^T \beta)
\]

where $g$ is the \textbf{link function} and $Y$ follows an exponential family.

\textbf{Examples}:

\begin{enumerate}
	\item \textbf{Linear Regression}: $Y \sim \mathcal{N}(\mu, \sigma^2)$, $g(\mu) = \mu$ (identity link)
	\item \textbf{Logistic Regression}: $Y \sim \text{Bernoulli}(p)$, $g(p) = \log\frac{p}{1-p}$ (logit link)
	\item \textbf{Poisson Regression}: $Y \sim \text{Poisson}(\lambda)$, $g(\lambda) = \log \lambda$ (log link)
\end{enumerate}

\textbf{Information Geometry of GLMs}:

\begin{itemize}
	\item The parameter space forms an exponential family manifold
	\item Maximum likelihood estimation corresponds to orthogonal projection
	\item The Fisher information matrix provides the natural metric
\end{itemize}

\subsubsection{8.4 Neural Networks and Deep Learning}

\textbf{Softmax Activation}:
The softmax function is the natural parameter to expectation parameter mapping for categorical distributions:
\[
p_i = \frac{e^{x^i}}{\sum_{j=1}^k e^{x^j}}
\]

\textbf{Cross-Entropy Loss}:
For classification, the cross-entropy loss is:
\[
L = -\sum_{i=1}^k y_i \log p_i
\]

This is the KL divergence between the true distribution and the predicted distribution.

\textbf{Natural Gradient}:
Using the Fisher information metric for parameter updates:
\[
\theta^{(t+1)} = \theta^{(t)} - \alpha \mathcal{I}^{-1}(\theta^{(t)}) \nabla_\theta L
\]

This provides more efficient optimization for neural networks.

\subsection{9. Connection to Aitchison Geometry}

\subsubsection{9.1 Compositional Data as Exponential Families}

The exponential/mixture duality is closely related to the Aitchison geometry of compositional data:

\textbf{Categorical Exponential Family}:
\[
p(\zeta = i) = \frac{e^{x^i}}{\sum_{j=1}^k e^{x^j}}
\]

\textbf{Connection to Aitchison Operations}:

\begin{itemize}
	\item \textbf{Aitchison addition}: $p \oplus q$ corresponds to adding natural parameters
	\item \textbf{Aitchison scalar multiplication}: $\lambda \odot p$ corresponds to scaling natural parameters
	\item \textbf{clr transform}: Maps from simplex to natural parameter space
\end{itemize}

\subsubsection{9.2 Geometric Structure}

\begin{enumerate}
	\item \textbf{Categorical distributions} form an exponential family
	\item \textbf{Simplex space} corresponds to mixture coordinates (expectation parameters)
	\item \textbf{clr-transformed space} corresponds to natural parameters (e-flat coordinates)
	\item \textbf{Weighted norms} reflect the Fisher information structure
\end{enumerate}

\subsubsection{9.3 Why Weighted Norms are Natural}

\textbf{Fisher Information Structure}:
The Fisher information matrix for categorical distributions is:
\[
\mathcal{I}_{ij} = \delta_{ij} p_i - p_i p_j
\]

This suggests natural weights $w_i = 1/p_i$ for different probability components.

\textbf{Orlicz Norms and f-Divergences}:
The connection between Orlicz norms and f-divergences provides a bridge between:

\begin{itemize}
	\item Functional analysis (norms on function spaces)
	\item Information theory (divergences between distributions)
	\item Compositional data analysis (geometry of the simplex)
\end{itemize}

\subsubsection{9.4 Practical Implications}

\textbf{Suitable Norms for Compositional Data}:

\begin{enumerate}
	\item \textbf{Weighted $L^p$ norms}: Account for natural variability structure
	\item \textbf{Orlicz norms}: Provide robustness and flexibility
	\item \textbf{Information-theoretic optimality}: Connect to maximum likelihood and minimum divergence principles
\end{enumerate}

\textbf{Applications}:

\begin{itemize}
	\item \textbf{Ecology}: Species abundance data with natural weighting
	\item \textbf{Economics}: Budget allocation with importance weighting
	\item \textbf{Genomics}: Gene expression data with regulatory importance
\end{itemize}

\subsection{10. Summary and Future Directions}

\subsubsection{10.1 Fundamental Insights}

The exponential/mixture duality reveals a fundamental geometric structure in probability theory:

\begin{enumerate}
	\item \textbf{Exponential families} are parameterized by natural parameters $x$ (e-coordinates)
	\item \textbf{Mixture families} are parameterized by expectation parameters $u$ (m-coordinates)
	\item These coordinates are related by the \textbf{Legendre transform} of the log-partition function
	\item \textbf{KL divergence} measures distances in this geometric structure
	\item \textbf{Maximum entropy} and \textbf{minimum divergence} principles have beautiful geometric interpretations
	\item The \textbf{Pythagorean theorem} holds in this information-geometric space
\end{enumerate}

\subsubsection{10.2 Connections to Other Areas}

\textbf{Statistical Mechanics}:

\begin{itemize}
	\item Exponential families correspond to Boltzmann distributions
	\item Log-partition function is the free energy
	\item Temperature parameter relates to the scale of natural parameters
\end{itemize}

\textbf{Optimal Transport}:

\begin{itemize}
	\item Wasserstein distances provide alternative geometric structure
	\item Connections through entropic regularization
	\item Sinkhorn algorithm exploits exponential family structure
\end{itemize}

\textbf{Information Theory}:

\begin{itemize}
	\item Rate-distortion theory uses exponential family structure
	\item Channel capacity problems involve mixture families
	\item Mutual information decompositions use orthogonal projections
\end{itemize}

\subsubsection{10.3 Open Problems and Future Research}

\textbf{Infinite-Dimensional Extensions}:

\begin{itemize}
	\item Exponential families in function spaces
	\item Gaussian processes as infinite-dimensional exponential families
	\item Variational inference in infinite dimensions
\end{itemize}

\textbf{Non-Euclidean Geometries}:

\begin{itemize}
	\item Hyperbolic and spherical information geometries
	\item Manifolds with non-constant curvature
	\item Applications to hierarchical and tree-like data
\end{itemize}

\textbf{Computational Aspects}:

\begin{itemize}
	\item Efficient algorithms for high-dimensional exponential families
	\item Scalable variational inference methods
	\item Hardware acceleration for information-geometric computations
\end{itemize}

\textbf{Applications to Modern ML}:

\begin{itemize}
	\item Generative adversarial networks through divergence minimization
	\item Transformer attention mechanisms as exponential families
	\item Federated learning with information-geometric privacy
\end{itemize}

\subsubsection{10.4 Conclusion}

This duality provides the foundation for understanding statistical inference, machine learning, and the geometry of probability distributions. The geometric viewpoint not only offers mathematical elegance but also leads to practical algorithms and deeper theoretical insights that continue to drive progress in statistics, machine learning, and information theory.

The beauty of information geometry lies in its ability to unify seemingly disparate concepts—from maximum entropy principles to neural network optimization—under a single geometric framework. This unified perspective continues to inspire new research directions and applications across diverse fields of science and engineering.


\end{document}